\documentclass[a4paper,10pt]{article}
\usepackage[utf8]{inputenc}
\usepackage[spanish]{babel}
\usepackage[margin=0.5cm]{geometry}
\usepackage{amsmath}
\usepackage{amsfonts}

\begin{document}
\pagestyle{empty}

\begin{minipage}[t][0.48\textheight][t]{0.48\textwidth}
  \centering \textbf{Evaluación Electrotecnia 2 - Tema A} \\
  \rule{\textwidth}{0.4pt}
  \begin{enumerate}
    \item Un conductor que se desplaza con una velocidad
      $v=3\frac{m}{s}$ perpendicular a un campo magnético de
      inducción $B=1T$, sabiendo que la longitud es $l=20cm$.
      Calcular la \textbf{fuerza electro motriz} $F.E.M$ inducida en
      el conductor. (1 pt)
    \item Un conductor de longitud $l=30cm$ se mueve
      perpendicularmente al campo con una velocidad de
      $v=4\frac{m}{s}$. Calcular el valor de la \textbf{inducción
      $B$} si la F.E.M. inducida es de $E_m=18V$. (1,5 pts)
    \item Una bobina de $N=400$ espiras es recorrida por una
      corriente $I=20A$ que da lugar a un flujo magnético de
      $\Phi=0,001Wb$. Calcular el \textbf{coeficiente de
      autoinducción $L$} de la bobina. (1 pt)
    \item Calcular la \textbf{fuerza electromotríz} inducida en una
      bobina de \( L = 0.005 \, \text{H} \) cuando la corriente que
      circula por ella varía a razón de \( \frac{\Delta I}{\Delta t}
      = 0.1\, \text{A/s} \). (1,5 pts)
    \item Un capacitor está formado por 2 placas paralelas de
      $s=1m^2$ si el dieléctrico es el mica y la distancia entre
      placas de 0.01mm, calcular su capacidad. (2 pts)
    \item Un condensador que se conecta a una tensión de $V=200V$,
      adquiere cada armadura una carga de $Q=6x10^{-9}C$. Calcular su
      capacidad. (1 pt)
    \item Un condensador tiene una capacidad de 50 $\mu$F y se le
      aplica una tensión de 12 V. ¿Cuál es la \textbf{carga}
      almacenada en el condensador? (2 pts)
  \end{enumerate}
\end{minipage}
\hfill
\begin{minipage}[t][0.48\textheight][t]{0.48\textwidth}
  \centering \textbf{Evaluación Electrotecnia 2 - Tema B} \\
  \rule{\textwidth}{0.4pt}
  \begin{enumerate}
    \item Un conductor rectilíneo se desplaza perpendicularmente a la
      dirección de un campo magnético de inducción $B=18000Gauss$,
      con una velocidad de $v=2\frac{m}{s}$. Calcular la
      \textbf{longitud del conductor} si está dentro del campo
      magnético si que induce en el una F.E.M de $E=14V$. (1,5 pts)
    \item Un conductor que se desplaza con una velocidad
      $v=10\frac{m}{s}$  en un campo magnético de inducción $B=2T$,
      sabiendo que la longitud es $l=10cm$. Calcular la
      \textbf{corriente} que circula por el conductor si se conecta a
      una resistencia de $R_L=10\Omega$. (1 pt)
    \item Un solenoide de $N=200$ espiras, longitud $l=40cm$ y
      diámetro $\varnothing=2cm$, tiene un núcleo de aire. Calcular
      el \textbf{coeficiente de autoinducción} $L$. (1,5 pts)
    \item Calcular el \textbf{flujo $\Phi$} que produce una bobina de
      $N=3000$ espiras por la que circulan $I=2A$  si el coeficiente
      de autoinducción es igual a $L=0,006H$. (1 pt)
    \item Un capacitor de 2 placas paralelas de $10cmx10cm$ si el
      dieléctrico es el vacío y la distancia entre placas de 0.1mm,
      calcular su capacidad. (2 pts)
    \item La carga que adquiere un condensador al conectarlo a
      $V=900V$ es de $Q=0,007C$. Calcular su capacidad. (1 pt)
    \item Un condensador almacena una carga de 0.25 C cuando se le
      aplica una tensión de 25 V. ¿Cuál es la \textbf{capacidad} del
      condensador? (2 pts)
  \end{enumerate}
\end{minipage}

\vfill

\begin{minipage}[t][0.48\textheight][t]{0.48\textwidth}
  \centering \textbf{Evaluación Electrotecnia 2 - Tema C} \\
  \rule{\textwidth}{0.4pt}
  \begin{enumerate}
    \item Calcular la \textbf{velocidad} que debe llevar un alambre
      de $l=20cm$ que se desplaza en un campo de inducción
      $B=1.4T$, para que la \textbf{F.E.M. inducida} sea de $E=2V$. (1,5 pts)
    \item Si un conductor tiene $l=0,5m$ y $B=2T$, se desplaza con
      una velocidad de $v=4\frac{m}{s}$.
      Calcular la \textbf{F.E.M} inducida $E$, la intensidad de la
      corriente $I$ conectando una resistencia de $R=12\Omega$. (2 pts)
    \item Calcular el \textbf{coeficiente de autoinducción $L$} de
      una espira $N=1$, si al circular por ella una corriente de
      $I=1A$ da lugar a un flujo en la sección transversal de la
      misma de $\Phi=200000Mx$. (1 pt)
    \item Un solenoide de $N=4000$ espiras y coeficiente de
      autoinducción igual a $L=0,01H$ debe producir por efecto de
      autoinducción un flujo de $\Phi=1000Mx$. Calcular la
      \textbf{corriente I}. (1,5 pts)
    \item Se desea saber la capacidad de el capacitor que posee
      placas paralelas de $S=0,1m^2$ si el dieléctrico es el vacío
      $\epsilon_0$ y la distancia entre placas de 0.001mm. ¿Cómo
      cambia el valor de la capacidad si el dieléctrico es
      \emph{vidrio}? (2 pts)
    \item ¿A qué tensión debemos conectar $C=104F$
      para que $Q=2x10^{-6}C$? (1 pt)
    \item Un condensador con una capacidad de 100 nF se carga hasta
      10 V. ¿Cuál es la \textbf{carga} en el condensador? (1 pt)
  \end{enumerate}
\end{minipage}
\hfill
\begin{minipage}[t][0.48\textheight][t]{0.48\textwidth}
  \centering \textbf{Evaluación Electrotecnia 2 - Tema D} \\
  \rule{\textwidth}{0.4pt}
  \begin{enumerate}
    \item Un conductor metálico de longitud $l=0.8 m$ se desplaza a
      través de un campo magnético uniforme con una velocidad de $v=3
      m/s$. La fuerza electromotriz (fem) inducida en el conductor es
      de $E_m=0.6 V$. Calcula la \textbf{inducción} ($B$). (1,5 pts)
    \item Un conductor metálico se desplaza a través de un campo
      magnético uniforme con una velocidad de $v=4 m/s$. La fuerza
      electromotriz (fem) inducida en el conductor es de $F.E.M=1,2
      V$. Si la intensidad del campo magnético es de $B=0,2 T$, ¿cuál
      es la \textbf{longitud} del conductor dentro del campo
      magnético en centímetros? (2 pts)
    \item Una bobina de autoinducción \( L = 0.01 \, \text{H} \)
      tiene una corriente que varía a una tasa de \( \frac{di}{dt} =
      -0.5 \, \text{A/s} \). ¿Cuál es la \textbf{fuerza
      electromotriz} inducida? (1,5 pts)
    \item Se tiene una bobina con un coeficiente de autoinducción $L$
      de $2$ H. Si el número de vueltas $N$ es de $1000$, el radio
      $r$ es de $0.05$ m y la longitud $l$ es de $0.1$ m, ¿cuál es la
      permeabilidad magnética $\mu$ del núcleo de la bobina? (2 pts)
    \item Calcular la energía que tiene un capacitor de $10\mu F$ si
      este tiene una carga de $Q=2C$. (1 pt)
    \item calcular la carga de un condensador de $C=200pF$ si se
      conecta a $V=100V$. (1 pt)
    \item Un condensador se carga hasta 200 mC cuando se le aplica
      una tensión de 400 V. ¿Cuál es la \textbf{capacidad} del
      condensador? (1 pt)
  \end{enumerate}
\end{minipage}

\end{document}
